\chapter{Introduzione} % senza numerazione
\label{cha:introduzione}

Nell'attuale era digitale, i dati rappresentano la risorsa strategica
centrale per ogni tipo di organizzazione, dalla sanità al commercio, fino
allo sport professionistico. Nonostante i database relazionali (RDBMS)
siano il motore dinamico della ricerca scientifica e del business, il
loro accesso è ancora limitato dalla barriera tecnica del linguaggio
SQL.

\section{Presentazione del
problema}\label{presentazione-del-problema}

La necessità di una sintassi rigorosa rende i dati accessibili quasi
esclusivamente a figure specializzate, costringendo gli utenti non
tecnici (manager, medici o analisti) a dipendere dai reparti IT. Questa
intermediazione genera inefficienze critiche: rallentamenti decisionali,
elevati costi di formazione e il rischio di incomprensioni tra chi pone
la domanda e chi scrive il codice, lasciando spesso inespresso il reale
potenziale del patrimonio informativo aziendale. La spinta principale di
questo lavoro risiede nella democratizzazione dell'accesso ai dati. In
un mercato che richiede risposte immediate, è essenziale trasformare i
database da archivi statici in strumenti di conversazione aperti a
tutti. Automatizzare la traduzione in SQL significa garantire autonomia
decisionale e migliorare l'efficienza operativa: il valore di
un'intuizione non deve più essere limitato dalla conoscenza del codice,
ma solo dalla capacità umana di porre le domande giuste. Il settore del
Text-to-SQL è stato rivoluzionato dai Large Language Models (LLM).
Sebbene modelli generalisti come GPT-4 abbiano aperto la strada,
l'attenzione si è oggi spostata su soluzioni specializzate come
Qwen-Coder o modelli fine-tuned come SQLCoder, capaci di gestire schemi
complessi con maggiore precisione. Parallelamente, l'introduzione della
tecnica RAG (Retrieval-Augmented Generation) ha segnato un punto di
svolta, permettendo al sistema di ``recuperare'' metadati e la
documentazione per costruire prompt contestualizzati. Questo approccio
riduce le allucinazioni dell'IA e permette ai modelli di adattarsi a
database specifici senza necessità di riaddestramento. Tale evoluzione
ha favorito la nascita di strumenti pronti all'uso come Text2SQL o
AskYourDatabase, che iniziano a integrare queste tecnologie con
interfacce intuitive per l'utente finale.

\section{Obiettivi e Risultati
Attesi}\label{obiettivi-e-risultati-attesi}

L'obiettivo principale di questo lavoro di tesi risiede nello sviluppo,
nell'implementazione e nella validazione sperimentale di un sistema
modulare a oggetti in linguaggio Python. Tale sistema ingegnerizzato deve tradurre in
modo efficace e sintatticamente coerente richieste formulate in
linguaggio naturale in interrogazioni SQL direttamente eseguibili su un
DBMS target. Il nucleo del progetto si focalizza sulla creazione di uno
strato software strutturato che opera come ponte tra l'utente (o
un'applicazione client) e la complessità intrinseca degli schemi
relazionali, garantendo un'elevata precisione nella mappatura tra i
concetti espressi verbalmente e le strutture logiche delle tabelle.

L'intento è fornire uno strumento flessibile, disaccoppiato e facilmente
integrabile in ecosistemi software preesistenti, abilitando piattaforme
terze a implementare funzionalità di analisi dati in modalità
self-service. Attraverso l'impiego di modelli di linguaggio avanzati
(LLM) e di tecniche evolute di gestione del contesto, l'obiettivo finale
dimostra la fattibilità dell'abbattimento della barriera tecnica del
codice, trasformando l'interrogazione dei dati in un processo immediato,
sicuro e accessibile.

Per conseguire questo traguardo, il lavoro di ricerca e sviluppo si
articola in una serie di obiettivi intermedi focalizzati
sull'automazione del workflow di traduzione, esecuzione e raffinamento:

\begin{itemize}
\item
  \textbf{Estrazione e normalizzazione dello schema:} Automatizzare la
  mappatura del database target sia tramite connessione diretta e
  recupero dei metadati dizionariali, sia attraverso l'induzione e la
  normalizzazione di input non strutturati (istruzioni DDL) in un
  formato JSON standardizzato, successivamente indicizzato all'interno
  di un database vettoriale.
\item
  \textbf{Ottimizzazione del contesto:} Ridurre la finestra di calcolo
  (context window) dei modelli linguistici estraendo
  dinamicamente, tramite euristiche di recupero vettoriale, solo le
  entità strettamente necessarie all'intento espresso dall'utente,
  arricchendo lo schema con regole di business e metadati
  extrastrutturali.
\item
  \textbf{Generazione tramite prompt dinamici:} Sviluppare un motore di
  compilazione dei prompt che aggrega algoritmicamente schemi
  relazionali, percorsi di giunzione (join paths) e lo storico
  dei fallimenti passati, garantendo il completo disaccoppiamento logico
  dai modelli linguistici sottostanti.
\item
  \textbf{Validazione e raffinamento iterativo (Refinement):}
  Implementare un ciclo di feedback a ciclo chiuso che verifica a
  runtime la conformità sintattica ed esecutiva della query prodotta,
  intercettando le eccezioni sollevate dal database per guidare un
  processo autonomo di auto-correzione.
\item
  \textbf{Esposizione architetturale tramite API:} Ingegnerizzare
  l'applicazione separando nettamente la logica di backend dal frontend,
  esponendo le funzionalità di generazione query e di supervisione
  critica attraverso endpoint programmabili e facilmente interrogabili.
\end{itemize}

Dal punto di vista prettamente ingegneristico, il progetto si
concretizza nella realizzazione di un modulo API robusto e scalabile,
distribuito sotto forma di un pacchetto Python documentato e pronto
all'uso. Tale soluzione è concepita per essere integrata agevolmente in
applicazioni esterne, quali dashboard di Business Intelligence o agenti
conversazionali aziendali, dimostrando versatilità nel gestire schemi
relazionali eterogenei.

Un risultato chiave del framework è rappresentato dall'ottimizzazione
del contesto tramite metodologie RAG (\emph{Retrieval-Augmented
Generation}). L'efficacia del sistema di recupero e l'integrazione
dinamica dello schema nel prompt consentono al modello di operare con
successo anche su database non inclusi nella fase di addestramento
iniziale, puntando ad ottenere una notevole capacità di generalizzazione in
modalità zero-shot.

A completamento dello sviluppo, l'intera impalcatura software è
sottoposta a una massiccia campagna di test basata sui principali
benchmark scientifici del settore, quali Spider e BIRD. Attraverso di essi
sono stati ottenuti risultati quantitativi e qualitativi che rivelano
quali dei requisiti sono stati raggiunti dal framework. 
Il principale traguardo consiste nel raggiungimento di
un'elevata accuratezza nella traduzione (Execution Accuracy),
minimizzando i fenomeni di allucinazione dei modelli e garantendo che il
codice SQL prodotto sia pronto per l'esecuzione sul database target
senza interventi manuali. Di conseguenza, il sistema persegue una netta
riduzione della barriera tecnica nell'accesso ai dati, semplificando i
processi di analisi e permettendo a utenti non esperti di estrarre
informazioni complesse in tempo reale, diminuendo la dipendenza da
intermediari tecnici.

\section{Struttura della Tesi}\label{struttura-della-tesi}

Il presente elaborato è organizzato in capitoli, strutturati secondo un
percorso logico che spazia dall'inquadramento teorico fino alla
validazione empirica sul campo:

\begin{itemize}
\item
  \textbf{Capitolo 2 --- Stato dell'arte e Background:} descrive
  l'evoluzione tecnologica dei Large Language Models e analizza lo stato
  dell'arte delle metodologie Text-to-SQL e dei database vettoriali. Il
  capitolo include una rassegna critica dei principali strumenti
  commerciali attualmente disponibili sul mercato, evidenziandone i
  limiti strutturali a cui il presente lavoro intende rispondere.
\item
  \textbf{Capitolo 3 --- Soluzione proposta:} illustra nel dettaglio
  l'architettura ingegnerizzata a oggetti, descrivendo le responsabilità
  dello \texttt{SchemaOrchestrator} e del \texttt{QueryOrchestrator}.
  Vengono analizzati i flussi logici del back-end FastAPI e
  l'interfaccia di front-end in React concepita per l'interazione con
  l'utente finalizzandone i casi d'uso.
\item
  \textbf{Capitolo 4 --- Test e valutazioni:} espone la metodologia
  sperimentale adottata per la validazione del framework sui benchmark
  internazionali Spider e BIRD. In questa sezione vengono analizzati in
  modo esaustivo i risultati quantitativi ottenuti su molteplici modelli
  ed ecosistemi di database, arricchendo la trattazione con uno studio
  inferenziale e statistico delle correlazioni.
\item
  \textbf{Capitolo 5 --- Conclusione:} riassume i principali traguardi
  conseguiti in relazione ai requisiti progettuali e delinea le
  promettenti traiettorie di ricerca future volte all'ottimizzazione del
  sistema.
\end{itemize}
